% =============================================================================
%  RL-H7-01 | Integrasi laporan Sprint 1
% -----------------------------------------------------------------------------
%  Role    : RL — Research Lead
%  Sprint  : H7 | 09.00–13.00 | SEKUENSIAL (menunggu tulisan semua peran)
%  Ticket  : (isi dari papan kerja)          Status: STARTER ONLY
%  Target  : reports/sprint1/sprint1_report.tex -> .pdf
%  Input   : 01_pendahuluan, 02_related_work, 03_metode (RL),
%            04_data (MDE), 05_iklim_dan_regime (CGS), 06_eksperimen_dan_hasil (MLE)
%  Output  : laporan Sprint 1 + lampiran run_id
%  Handoff : ALL-H7-01 (Gate 3)
%  DoD     : [ ] setiap angka punya rujukan run_id (footnote/tabel lampiran)
%            [ ] hasil negatif ditulis eksplisit
%            [ ] MDE sudah QA semua angka
%            [ ] batas klaim diperbarui berdasarkan hasil
% =============================================================================
\documentclass[11pt,a4paper]{article}
% =============================================================================
%  00_preamble.tex | shared preamble untuk semua deliverable RL
%  Job     : bukan job card; dipakai oleh RL-H1-01 .. RL-H7-01
%  Target  : docs/ (salin bersama file .tex)
%  Catatan : kompilasi dengan pdflatex. Untuk XeLaTeX/LuaLaTeX, hapus inputenc.
% =============================================================================

\usepackage[utf8]{inputenc}
\usepackage[T1]{fontenc}
\usepackage{lmodern}
\usepackage[a4paper,margin=2.4cm]{geometry}
\usepackage{amsmath,amssymb}
\usepackage{booktabs}
\usepackage{longtable}
\usepackage{array}
\usepackage{graphicx}
\usepackage{xcolor}
\usepackage{enumitem}
\usepackage{float}
\usepackage{caption}
\usepackage{url}
\usepackage{microtype}
\usepackage[colorlinks=true,linkcolor=blue!50!black,urlcolor=blue!50!black,citecolor=blue!50!black]{hyperref}

% --- helper wajib -----------------------------------------------------------
% Path/nama file aman (underscore, garis miring) di dalam teks.
\newcommand{\fpath}[1]{\texttt{\detokenize{#1}}}
% TODO yang mencolok. Jangan dihapus sebelum job selesai.
\newcommand{\todo}[1]{\textcolor{red!75!black}{\textbf{[TODO: #1]}}}
\newcommand{\jobid}[1]{\textsf{\textbf{#1}}}
\newcommand{\runid}[1]{\texttt{\detokenize{#1}}}
\newcommand{\claim}[1]{\textsf{#1}}
\newcommand{\good}{\textcolor{green!45!black}{$\checkmark$}}
\newcommand{\bad}{\textcolor{red!75!black}{$\times$}}
\newcommand{\maybe}{\textcolor{orange!80!black}{$\sim$}}

\newcolumntype{L}[1]{>{\raggedright\arraybackslash}p{#1}}
\newcolumntype{C}[1]{>{\centering\arraybackslash}p{#1}}

\setlength{\parindent}{0pt}
\setlength{\parskip}{4pt}
\captionsetup{font=small,labelfont=bf}
\setlist[itemize]{leftmargin=1.2em,itemsep=2pt,topsep=2pt}
\setlist[enumerate]{leftmargin=1.5em,itemsep=2pt,topsep=2pt}

% Judul job yang konsisten di semua dokumen.
\newcommand{\jobheader}[4]{%
  \begin{center}
  {\Large\bfseries #1}\\[2pt]
  {\large #2}\\[4pt]
  \begin{tabular}{@{}ll@{}}
  \toprule
  Role & #3 \\
  Sprint & #4 \\
  \bottomrule
  \end{tabular}
  \end{center}
}


\title{\jobid{RL-H7-01} | Laporan Sprint 1\\
\large Robust Climate-Conditioned Dynamic Graph Learning for\\
Food Price Forecasting Under Extreme Weather and Climate Regime Shifts}
\author{RENARO Lab \\ \small RL, MDE, MLE, CGS}
\date{23 November 2026 \quad | \quad \textcolor{red!75!black}{DRAFT UNTUK GATE 3}}

\begin{document}
\maketitle

\begin{table}[H]
\centering
\small
\begin{tabular}{@{}L{3.6cm}L{10.2cm}@{}}
\toprule
Versi & \todo{v0.9 draft / v1.0 final} \\
Commit & \todo{\texttt{sha}} \\
Dataset & \texttt{dataset\_v1} | hash \todo{} \\
Protokol & \texttt{protocol-v1} | hash \todo{} \\
Tag model & \texttt{gate2-freeze} | hash \todo{} \\
Holdout dibuka & \todo{waktu} | \fpath{docs/holdout_seal.md} \\
\bottomrule
\end{tabular}
\end{table}

\section*{Ringkasan eksekutif}
\todo{Setengah halaman. Jawab RQ0--RQ3 masing-masing 1 kalimat + keputusan
Gate 1/2/3. Tulis angka utama dengan run\_id. Jangan melebihkan klaim.}

\begin{table}[H]
\centering
\small
\begin{tabular}{@{}L{1.4cm}L{8.8cm}L{4.0cm}@{}}
\toprule
RQ & Jawaban 1 kalimat & Bukti utama \\
\midrule
RQ0 & \todo{} & \todo{} \\
RQ1 & \todo{} & \todo{} \\
RQ2 & \todo{} & \todo{} \\
RQ3 & \todo{} & \todo{} \\
\bottomrule
\end{tabular}
\end{table}

\section{Pendahuluan}
\todo{Ringkas dari RL-H4-01: masalah, gap, RQ0--RQ3, kontribusi berurutan
benchmark $\to$ bukti empiris $\to$ metode, scope A, batas non-kausal.}

\section{Related work}
\todo{Ringkas 4 kelompok + 3 kalimat posisi dari \fpath{docs/positioning.pdf}.
Sitasi dengan \texttt{\textbackslash cite\{...\}}.}

\section{Data}
\todo{Ringkas dari MDE-H6-01: sumber, cakupan, cleaning, kalender, mask BPS,
graph prior, aturan availability, keterbatasan. Rujuk \fpath{docs/data_card.pdf}.}

\section{Iklim dan regime}
\todo{Ringkas dari CGS: CHIRPS + pembobotan produksi, indeks pentad, RONI vintage,
definisi regime, jumlah node-hari per regime. Tegaskan CHIRPS harian = disagregasi pentad.}

\section{Metode}
\todo{Ringkas Modul 1 (persamaan \ref{eq:module1} di RL-H4-01), Modul 2--3,
Modul 4 stretch, baseline, protokol. Jangan menyalin seluruh dokumen protokol.}

\section{Eksperimen dan hasil}
\subsection{Validasi (O1--O3)}
\begin{table}[H]
\centering
\small
\begin{tabular}{@{}L{3.2cm}C{1.6cm}C{1.6cm}C{1.6cm}C{1.8cm}C{1.8cm}@{}}
\toprule
Model & MAE & RMSE & sMAPE & Worst regime & DM $p$ \\
\midrule
Seasonal naive & \todo{} & \todo{} & \todo{} & \todo{} & -- \\
LightGBM (a/b/c) & \todo{} & \todo{} & \todo{} & \todo{} & \todo{} \\
GWN & \todo{} & \todo{} & \todo{} & \todo{} & \todo{} \\
MTGNN & \todo{} & \todo{} & \todo{} & \todo{} & \todo{} \\
Latent regime (RCSGNet-like) & \todo{} & \todo{} & \todo{} & \todo{} & \todo{} \\
Lag graph (MillGNN-like) & \todo{} & \todo{} & \todo{} & \todo{} & \todo{} \\
Module 1 & \todo{} & \todo{} & \todo{} & \todo{} & \todo{} \\
Module 1--3 & \todo{} & \todo{} & \todo{} & \todo{} & \todo{} \\
\bottomrule
\end{tabular}
\caption{Rata-rata 5 seed $\pm$ std; rujuk \runid{run\_id} di lampiran.}
\end{table}

\subsection{Uji mekanisme sintetis}
\todo{Ringkas \fpath{reports/synth/recovery_summary.md}: lag accuracy, Spearman,
delta lag, kontrol. Sebutkan jika gagal.}

\subsection{Robustness (RQ3)}
\todo{Ringkas \fpath{reports/robustness/robustness_summary.md}: kurva MAE vs p,
joint shift, per regime.}

\subsection{Holdout 2026}
\todo{Ringkas \fpath{reports/holdout/findings.md}: holdout hanya dibuka sekali;
model di luar manifest tidak dijalankan; hasil negatif dilaporkan.}

\section{Batas klaim (diperbarui)}
\begin{table}[H]
\centering
\small
\begin{tabular}{@{}L{2.2cm}L{8.2cm}L{4.2cm}@{}}
\toprule
Kategori & Klaim & Status bukti \\
\midrule
AMAN & \todo{} & \todo{} \\
BERISIKO & \todo{} & \todo{} \\
JANGAN & graph adaptif/gating/invariance sebagai novelty; edge = arus/kausalitas;
``pertama memakai iklim pada GNN harga pertanian'' & tetap dilarang \\
\bottomrule
\end{tabular}
\end{table}

\section{Rekomendasi Sprint 2}
\todo{Pilih jalur A/B/C sesuai Gate 3 dan brief \S10.1. Tulis 5--8 item backlog:
metode, data, robustness, tulisan. Sebut sumber daya yang dibutuhkan.}

\appendix
\section{Lampiran: peta angka ke run}
\begingroup
\footnotesize
\setlength{\tabcolsep}{3pt}
\begin{longtable}{@{}L{4.0cm}L{2.8cm}L{1.8cm}C{1.2cm}L{3.8cm}@{}}
\toprule
Angka/tabel & run\_id & Model & Seed & Catatan \\
\midrule
\endhead
\todo{} & \todo{} & \todo{} & -- & \todo{} \\
\bottomrule
\end{longtable}
\endgroup

\section{Lampiran: artefak}
\begin{table}[H]
\centering
\small
\begin{tabular}{@{}L{5.0cm}L{9.0cm}@{}}
\toprule
Artefak & Path/hash \\
\midrule
Dataset v1 & \fpath{data/processed/dataset_v1/} | \todo{} \\
Protokol v1 & \fpath{docs/eval_protocol_v1.pdf} | \todo{} \\
Manifest beku & \fpath{frozen/manifest.yaml} | \todo{} \\
Prediksi holdout & \fpath{reports/holdout/predictions/} | \todo{} \\
Kode rilis & \fpath{release/code/} | \todo{} \\
\bottomrule
\end{tabular}
\end{table}

\bibliographystyle{plain}
\bibliography{references}

\end{document}
